\documentclass{article}
\usepackage[margin=1in]{geometry}
\usepackage{amsmath,amssymb}
\usepackage[UTF8]{ctex}
\begin{document}
\section*{随机取货和送货，结合预测性路线规划}
\paragraph{条件性源解释。}
以下模型选择了2015年论文中的(4)-(6)节中的路线-状态维度、普通路段描述以及出发/装载定义（印刷于\ 241--242页），并在图\ 1中使用了物理当前位置锚点。该模型将接驳指标作为等待因素。2015年的文章指出，该模型的构建方式总结了Cort\'等人\ （2009年）的研究；该论文的印刷版\ 33中明确将等待因素与已计算成本的路线接驳指标相对应。2008年的早期论文使用$z$作为这一指标的正式名称。这里的旅行和等待系数遵循2015年的两种术语（8）。下面具体说明了所选表达式与印刷矩阵之间的差异。

\paragraph{信息、历史与决策。}
考虑一个已观察到的当前请求，以及由$N-1$构成的未来请求树，这些请求都带有历史转移概率。 有$F$辆车辆。
在历史记录$h$中，深度为$d\in\{0,\ldots,N-1\}$时，让$e_h$成为已知或预测请求的时间点；同时，让$\pi_h$成为达到$h$的概率。 该根的概率为一。 所有具有相同观察结果的历史记录，在分支点之前，都会做出相同的决策。
每辆车辆的初始位置、装载情况、前方的路线以及未被服务的乘客信息都已被明确给出。 为了生成儿童历史中的状态，需要按照父母指定的路径进行，直到达到儿童请求的时间点。 移除已完成的前缀，包括服务在请求时间结束的情况，保留所有仍在船上的乘客，并插值车辆在其当前行驶段的位置。 闲置的车辆会保持在其最后的位置。
在十个输入案例中，旅行速度始终保持不变，且每个记录的服务和转移时间均为零。

在每一段历史记录中，为即将到来的乘客选择一辆车辆$j$，一个取货等级$a$以及一个送货等级$b$。在插入数据之前，先确定该车辆的待处理停靠点$q$，然后填入$w=q+2$。序数域分别为$1\leq a\leq w-1$、$2\leq b\leq w$和$a<b$。从新列表中删除新增的取货和送货信息，以恢复原有的有序列表。其他车辆则保持其原有的待处理路线。要求取货必须在送货之前完成，同时要求乘客与车辆的一致性、车辆上的乘客数量以及车辆容量都符合规定。初始状态下的乘客在车辆行驶完其路线后，其待处理的送货和卸客任务也一并完成。在当前的时间范围内执行根决策。

\paragraph{已发布的路线的状态。}
对于车辆$j$和位于$h$的出发路线，普通停靠点为$i=1,\ldots,w_{jh}$，而零行则是该车辆的锚点。将$r_{jhi}=1$填写为取货点，将$r_{jhi}=0$填写为交货点。该停靠点还记录了乘客的身份以及从前一个点的行驶时间$\Gamma_{jhi}$，其中包括换乘的延迟时间。在指定的历史数据中，条件状态预测结果如(6)所示。
\[
 L_{jh0}=L^{\rm anchor}_{jh},\quad T_{jh0}=e_h,\qquad
 L_{jhi}=L_{jh,i-1}+2r_{jhi}-1,\qquad
 T_{jhi}=T_{jh,i-1}+\Gamma_{jhi}.
\]
这里的$L_{jhi}$指的是离开停止点$i$之后的负载；前一个停止点的负载则沿着路径$i$继续移动。向量宽度为$w_{jh}+1$。根据输入的行驶时间和服务时间以及固定的路线，条件期望值等于这些数值。在后续的操作中，会生成一个新的状态，该状态将决定后续的路线。

\paragraph{成本。}
使用同一路线上的 pickup 因子，该路线的等级和状态正在被定价。所涉及的路线组$S_h$的舰队成本为
\[
 D_h=\sum_{j=1}^{F}\sum_{i=1}^{w_{jh}}
 \left[L_{jh,i-1}\left(T_{jhi}-T_{jh,i-1}\right)
       +r_{jhi}\left(T_{jhi}-T_{jh0}\right)\right].
\]
以相同的规则对前述路线进行定价，从而得到$D_{-1}$。将$p(h)$定义为$h$的父节点，而$D_{p(h)}=D_{-1}$则作为根节点。这个有限的决策树能够最小化(7)的期望值增量。
\[
 \min_{\text{feasible nonanticipative insertions}}
       \sum_{d=0}^{N-1}\sum_{h:\,\mathrm{depth}(h)=d}\pi_h\big(D_h-D_{p(h)}\big).
\]
孩子们的概率之和等于父母的概率，因此相同的目标在最后的决策历史中等于预期成本减去原始基线。对每辆车和每个合法等级对进行贝尔曼最小化处理，可以确定有限输入下的最优解。

\paragraph{印刷版教材。}
在 2015 年（4–6），预测的状态$x(k+t)$的宽度与已发布的$S(k+t-1)$相同。而（8）则分别描述了$w_j(k+t)$和$z_j^i(k+t)$，其中$t=N$处控制向量结束于$k+N-1$。上述选择符合排名状态及已发布列表的要求，并使用$k$至$k+N-1$作为决策时刻。2009年（11）的文献中使用了之前的已发布列表及其取货因素，以及预测的快照数据。而2015年（8）则没有提到驱动器的费用以及额外的等待权重。2015年的附录A显示了一个严格下三角的负载矩阵，以及$B_T^2=\operatorname{diag}(0,1,\ldots,1)$作为出发时间。对于带符号的停止变化$d_i=2r_i-1$，其负载在排名$i\geq1$时为$L_0+d_0+\sum_{s<i}d_s$。上述停止后的定义则给出$L_0+\sum_{s\leq i}d_s$。对角线上的时钟值为$T_0+\Gamma_i$，而累积时钟值为$T_0+\sum_{s\leq i}\Gamma_s$。原始的文字转录内容被单独保留下来以便比较。传统输入中还包含了乘客特定的价格系数；根据所选的指标解释方式，这些系数并没有实际的操作效果。每个输入都保留了其原始的历史树、尺寸、需求、概率以及车辆数据。
\end{document}
